% =============================================================================
%  Seccion: Metodologia
%  Contenido de EJEMPLO. Muestra ecuaciones numeradas y referenciadas.
% =============================================================================

\section{Metodologia}\label{sec:metodologia}

Que se hizo y por que ese metodo y no otro. Con suficiente detalle para que
alguien mas lo pueda reproducir.

Una formula corta va dentro del parrafo, entre signos de pesos: la matriz de
covarianzas muestral es $S = \frac{1}{n-1} X^{\top} X$ cuando $X$ ya esta
centrada.

Cuando la formula es el punto del parrafo, va aparte y numerada. La
descomposicion espectral de $S$ se escribe como

\begin{equation}\label{eq:espectral}
  S = P \Lambda P^{\top},
  \qquad
  \Lambda = \operatorname{diag}(\lambda_1, \dots, \lambda_p),
\end{equation}

donde $\lambda_1 \geq \dots \geq \lambda_p \geq 0$ son los valores propios y las
columnas de $P$ los vectores propios. La proporcion de varianza que explica la
$k$-esima componente se sigue de \cref{eq:espectral}:

\begin{equation}\label{eq:varianza}
  \pi_k = \frac{\lambda_k}{\sum_{j=1}^{p} \lambda_j}.
\end{equation}

Para varias ecuaciones alineadas se usa \verb|align|, y el \verb|&| marca por
donde se alinean (aqui, por el signo igual):

\begin{align}
  Z_1 &= p_{11} X_1 + p_{21} X_2 + \dots + p_{p1} X_p, \label{eq:pc1} \\
  Z_2 &= p_{12} X_1 + p_{22} X_2 + \dots + p_{p2} X_p. \label{eq:pc2}
\end{align}

Si una ecuacion no necesita numero, usa la version con asterisco
(\verb|equation*|, \verb|align*|) en lugar de dejar el numero colgado sin que
nadie lo cite.

El procedimiento completo sigue a \textcite{hastie2009} y se implemento en
\texttt{src/}.
